\section{Relatórios}

Eu gostaria de acrescentar uma rápida revisão baseada em dois relatórios globais sobre a desigualdade no mundo. Opto por este ponto de partida pois acredito que ele pode nos motivar a refletirmos sobre se este é um assunto pertinente. Me basearei aqui principalmente em dois relatórios. Vamos começar com o relatório da ONU intitulado ``Sustainable development outlook 2019 : gathering storms and silver linings''\cite{ONU}.  

\subsection{ONU} 

%https://github.com/jdansb/Econophysics/blob/main/files/Relatorio_desigualdade.md#user-content-fnref-1-1b8a500008960ecb392c37d88868ec69

The Sustainable Development Goals were conceived within the 2030 Agenda for Sustainable Development and are among the greatest hopes for a sustainable future, but they have faced significant challenges. In the words of the report’s preface:

\begin{quote}
Persistently high levels of inequality entrench uncertainty and insecurity among people, reinforcing divisions and undermining trust in institutions and government.
\end{quote}

The weakening of global economic growth and the rise of economic inequality pose a serious threat to the implementation of the Sustainable Development Goals. For these goals to be achieved, the international community must look beyond mere economic growth. An increase in GDP does not necessarily reflect an improvement in quality of life.

There remains a persistently high level of job insecurity in the world: most of the population is engaged in informal or vulnerable work, and one in three adults faces chronic risks of unemployment.

It is necessary to reduce inequality in all its forms and manifestations. The 2030 Agenda has embraced the motto “leave no one behind” as a guiding principle. However, in practice, there is little evidence that the world is winning the fight against inequality. This inequality prevents sustainable development in many ways: through economic stagnation, limits to social mobility and human development, and by undermining trust in institutions.

In practice, the effects of inequality can be seen, for example, in the fact that violent conflicts often arise from group-based inequalities and housing insecurity. It is therefore necessary to expand economic opportunities and ensure housing, especially for the most marginalized and vulnerable groups.

In developing countries, the population spends about 80\% of total income on basic necessities (food, housing, healthcare, and energy), while even in developed countries, spending on basic resources reaches nearly 60\%. One in five people in the world lives in inadequate housing, and one billion people still live in informal settlements. To address this situation, it is necessary to act across all dimensions by strengthening social protection measures and promoting the redistribution of wealth, income, and opportunities.

It is clear that the world has tolerated rising inequality for far too long and must now face its consequences — which erode solidarity and social cohesion. National and international efforts must be directed toward taming and reducing inequality, and this can only be achieved through regulatory measures. One thing is clear:*economic growth alone will not lift everyone out of poverty, and policies must look beyond GDP growth. TThe best results achieved so far in the pursuit of development come from a combination of national policies, such as those aimed at the labor market, and redistribution policies, which protect the most vulnerable and marginalized.

In most countries, wage growth has not kept pace with productivity, largely due to the absence of proactive policies. Wages have remained stagnant for the bottom 50\% of society, while they have increased substantially for the top 10\% and 1\%. Meanwhile, the cost of living has risen, disproportionately affecting the lower tiers of the social pyramid. Data collected in 2019 show that in the United States, 15\% of the population lives in high financial insecurity and 25\% in moderate insecurity. The report already pointed out that financial insecurity influences voting behavior toward the far right.

As for the rest of the world, more than two-thirds of countries are experiencing a rise in income and wealth inequality. There exists a reinforcing mechanism between vulnerable employment\footnote{Vulnerable workers are employed in low-skilled, low-paid jobs with little job security. They are often exploited by their employers and forced to work long hours for low pay, or to work in dangerous conditions.} and inequality levels: on one hand, high levels of vulnerable employment reinforce persistent inequality, while on the other, high inequality contributes to the expansion of vulnerable work. %O que é trabalhador \cite{vulneravel}

Given this scenario, the public sector cannot be excluded from the debate, since it is an observable fact that the concentration of wealth also results in a concentration of political power. In the United States, public institutions often act in complicity with policies that favor a small, wealthy fraction of the population at the expense of the majority. This high concentration of political power allows the rich to impose policies that serve their own interests, even when such policies do not benefit the broader public.

In this way, the concentration of power at the top of the social pyramid enables a decision-making process that disregards and fails to meet the needs of people outside this privileged segment. Seeking to protect their own interests, this minority fraction of the population—those who hold political and economic power—has both the incentive and the means to actively block the introduction of inequality-reducing measures, effectively creating an inequality trap.

For instance, higher inequality levels are associated with lower union density and weaker public labor protections.

When we examine global data, we can see that inequality between countries, as measured by the Gini coefficient, has decreased. However, this result is largely due to the massive populations of Asian countries whose average wealth and income per capita have historically been below the global mean. As countries such as China and India converge toward global averages, the global Gini coefficient declines mechanically.

By contrast, inequality within countries, also measured by the Gini index, has increased in nations that together account for over two-thirds of the world’s population. It is also necessary to interpret the Gini index critically: many countries that reported a decline in the Gini coefficient simultaneously experienced a rise in the share of wealth held by the top 10\%, since the Gini is less sensitive to variations at the extremes of the distribution.

The report further highlights that even countries showing a decline in inequality according to different indices have experienced wage growth and job creation that remain too slow to allow those at the bottom of the income spectrum to escape poverty.

Other consequences of inequality include its restrictive impact on social mobility. High inequality is closely linked to low social mobility, undermining the meritocratic narrative of capitalism, as poorer households face financial barriers that prevent them from investing adequately in their children’s education and opportunities for advancement.

Inequality also manifests in how financial crises affect different social classes unevenly. While the middle class suffers disproportionately from economic downturns, the upper class not only has greater capacity to prepare for crises but often even receives financial compensation from banks. Meanwhile, the lower class tends to remain geographically and economically isolated from both national and global markets due to its structural position within the economic system.

The globalization of trade has often been cited as a key factor in the rise of inequality within both developed and developing countries. However, this claim is not strongly supported by empirical evidence. Cross-country data from recent decades show only a limited correlation between trade openness and domestic inequality. According to recent studies, although trade liberalization plays a role in income inequality, its impact is secondary to that of technological development, financial openness, and financial depth. The main drivers are in fact poor public policies, such as the weakening of workers’ bargaining power, inadequate educational development, and the dismantling of social safety nets.

An alarming fact shows that 40\% of the population in OECD countries is economically vulnerable, meaning they lack the financial capital to live for three months without income above the poverty line. There are far more economically vulnerable people than those classified as income-poor. Situations of insecurity are directly linked to inequality levels: very high inequality intensifies feelings of insecurity and uncertainty. Conversely, growing insecurity has different impacts on different population groups, which can further aggravate inequality.

For instance, economic insecurity affects cognitive performance. Laboratory experiments show that people perform differently depending on the gap between what they need and the resources they have available to achieve it. The larger this gap, the more awareness of the situation captures attention and generates intrusive thoughts, impairing cognitive capacity. Those facing economic insecurity bear a kind of “cognitive tax”, which affects their well-being and productivity.

Inequality also impacts household resilience to climate events, which depends heavily on socioeconomic conditions. It undermines public trust in governments, institutions, and even among individuals, creating a political environment favorable to candidates considered outside the mainstream—especially populists and nationalists.

This divide between elites and regular workers at the national level has also been projected internationally, with nationalist forces capitalizing on domestic discontent and blaming multilateral institutions for serving elite interests. The multilateral trading system is in crisis; difficulties in addressing tax evasion, for example, reflect inadequate international regulation.

Unless major economic powers make serious efforts to reform the multilateral system, improve international cooperation on taxation, facilitate technology exchange, and address the climate crisis. less-developed countries will be left behind—leaving their most vulnerable populations to bear the brunt of the consequences.

Between the late 1990s and late 2000s, public support for redistribution increased in almost 70 per cent of the advanced and developing economies surveyed. However, this has not materialized due to the greater political power of the rich compared to the poor. The recent rise in inequality across the world is linked to the elimination of progressive tax schemes, the reduction of income taxes (especially at higher levels), and other policies that benefit the wealthiest segments of society. Among the proposed measures to fight inequality, the report identifies taxation as the main mechanism to be used. Progressive tax programs are an important factor in reducing inequality, while taxes on capital gains and inheritance also help to reduce intergenerational inequality.

Regarding tax targets, it is stated that reducing inequality in purchasing power rather than income is more effective in lowering consumption inequality. Once money is collected through tax programs, it must be redistributed according to public policy. Redistributive social policies may include cash benefits, such as subsidies, as well as non-monetary benefits, such as public health and education programs. Latin American countries have successfully experimented with cash transfer programs, and proposals for a universal basic income have emerged in several parts of the world. Land reforms also have a history of helping to fight inequality.

Although it is not yet clear what this structural change should look like, it is known that a proactive approach from states and lawmakers is necessary, since leaving these ambitions to the market will not lead to the required structural transformation. Producers focus on immediate success and profit, often ignoring social and environmental costs. Consumers are constrained to buying goods and services at affordable prices. The current system of production and consumption is driven by mechanisms that fail to account for pollution and social costs. Awareness campaigns are slow and inadequate for the scale and urgency of environmental issues. There is no magic solution — what must be discussed are public policies to address climate change and inequality.

The report proposes shifting the discourse from praising the "fastest-growing economy" to praising the "country that reduced inequality the most." For this to be possible, structural transformations are not only acceptable but are an integral part of the strategy to address inequality. It is clear that public policies adopted in the past have been lamentably insufficient to combat poverty and achieve sustainable development; the best that can be said is that there has been some stability in the last 15 years, following deterioration in the previous 15 years.

Among the taxes suggested by the report is a carbon tax, aimed at creating a financial disincentive for CO2 emissions. The report also mentions property and income taxes. For much of the 20th century, countries implemented progressive taxes (higher taxes for higher-income individuals) to ensure more equitable redistribution and the provision of quality public services. However, based on the economic argument that lower taxes encourage private investment and generate a larger multiplier effect than public investment, and also due to the political power that high-income individuals can exert—propagating this argument—tax progressivity has declined in most developed economies over the past four decades, especially at the top.

Property taxes do not directly mitigate income or wealth inequality but are the main source of revenue to fund education in the U.S. This type of property tax is widely used in developed countries and less so in developing countries. Wealth taxes, which apply to real and financial assets, are becoming increasingly popular. But critics argue that it yields little revenue relative to administrative costs and can be easily circumvented through deductions, hiding funds, or relocating assets to other countries. This has led to the elimination of wealth taxes in many contexts. As of 2017, only seven OECD countries maintained such a tax.

The land tax is considered the most efficient of all taxes. However, it remains a novelty in most countries. The concept is simple: a significant portion of wealth is captured in the value of urban land, and this value increases primarily due to public investment in the area. Individual actions by the landowner do not increase the value. The land tax therefore targets the increase in land value attributable to public investment, an "unearned" gain for the landowner.

As a wealth redistribution measure, there is the proposal of a universal basic income (UBI). This is an idea that has existed at least since the late 18th century and has seen a resurgence in the 21st century: a program that provides every individual, regardless of any criteria, with enough money to cover basic needs. Proponents argue that it can allow individuals to refuse undesirable work, increasing workers' bargaining power and reducing inequality.

Opponents, however, argue that there is a negative incentive, functioning as a disincentive to work. It is important to note that the effectiveness of any public policy depends on how the program interacts with a broader welfare system. As a complement, UBI can have an equalizing effect; however, as a substitute for all programs, it could worsen inequality. The greatest challenge concerns the resources required to fund the program. Finland conducted a pilot study between 2017 and 2018, and preliminary results indicate that UBI had no effect on employment levels but increased confidence in individuals' physical and mental well-being, as well as trust in institutions.
Global Wealth Report

\subsection{UBS}
Complementando ao relatório da ONU, consciente que alguns leitores poderiam alegar um certo viés no relatório da ONU, eu procurei também alum relatório feito por alguma instituição privada. O que trago aqui são algumas informações retiradas do relatório global de riqueza  realizado pelo tradicional banco suiço Credit Suisse no mesmo ano\cite{UBS}.

The UN report is already quite comprehensive for an introduction. For a deeper dive into the topic, I would recommend Piketty's book Capital in the Twenty-First Century. However, since I already have some notes on the UBS report, I will use them to make a few additional comments. This report was created by two banks and is the most up-to-date source on wealth distribution in the world; interestingly, one of the two banks self-declares as the largest global wealth manager.

Much of the report's analysis seems impractical for my purposes because it places significant weight on the exchange rate of the dollar relative to other currencies. For example, one of the results presented is that global wealth decreased, but much of this is due to the dollar’s relationship with other currencies, as reported. Nevertheless, there are some data points worth highlighting. First, we see confirmation of the previous report that, to a large extent, the reduction in inequality between countries is driven by growth in emerging economies, especially China, but also largely India. There are also figures indicating relative stability in both the Gini coefficient and the wealth of the top 1\% since the 2000s.

Regarding wealth distribution among different fractions of the population in the United States, the median wealth of Black households is just over 10\% of the median wealth of White households. The median wealth of Hispanic households relative to White households is 20\%. As previously reported, crises affect different sectors of society differently. Part of the explanation for this phenomenon lies in the different nature of wealth across classes. For example, most of the Hispanic group has their investment in housing (physical or non-financial assets), whereas White men, who have the highest average and median wealth, hold most of their wealth in non-physical (financial) assets\footnote{Financial and non-financial assets differ based on how they are bought and sold. Many financial assets, such as stocks and bonds, are traded on exchanges and can be bought or sold on any business day the exchange is open. It is easy to obtain the current market price to buy or sell these assets. As long as the market is liquid, there will be a buyer for every seller and vice versa. On the other hand, a non-financial asset, such as equipment or a vehicle, can be difficult to sell because there is no active market of buyers and sellers. The price of the non-financial item may be uncertain, as there is no market standard. Instead, many non-financial assets are sold when the seller finds a potential buyer and negotiates a sale price. The time required to find a buyer, complete the sale, and transfer the physical asset makes non-financial assets illiquid.}. %\cite{investopedia1,investopedia2}

Regarding global numbers, the poorest 52.5\% of the population holds 1.2\% of all wealth. In the poorest countries, 80\% of the population is concentrated at the base of this pyramid, while in richer countries it is 30\% and more transient compared to poorer countries. This fraction of the global population also shows great diversity across the world. The top 1.1\% richest hold 45.8\% of wealth and are particularly concentrated in certain regions and countries, sharing a similar lifestyle and opportunities, with a focus on non-physical assets. This global wealth distribution indicates the existence of an international pyramid among countries as well.

In terms of figures, one needs \$9,000 to be among the top 50\% richest, 16 times that to be among the top 10\% richest (\$140,000), and 125 times that to be among the top 1\% richest (\$1,000,000). Returning to the issue of the 'decline' in global inequality, people with wealth close to the global median tend to live in emerging economies and have their wealth mainly in the form of physical assets, specifically housing. During the pandemic, the poorest population was economically impacted, needing to use their savings and incur debt. For the poorest population, non-financial assets in their possession are more than double the wealth they hold, but the value of these assets is offset by the debts they have. The richest population, on the other hand, was not only less vulnerable but actually benefited from the crisis.

The report also aims to reduce the issue of inequality to two questions:
\begin{itemize}
\item    How far is the top of the pyramid from the average citizen? This is frequently answered by the percentage of wealth held by the top 1% (or 10%).
\item    How far below the base of the pyramid is from the average citizen?
\end{itemize}
Looking at numbers from different countries, the top 1\% richest tend to hold between 30\% and 35\% of wealth, while the top 10\% hold between 60\% and 65\% of total wealth. Comparing countries, a Gini of 70 is generally seen as relatively low, and 80 as relatively high. Recall that the Gini coefficient measures inequality and ranges from 0 (total equality) to 100 (total inequality). Inequality increased in liberalizing countries, and according to this report, decreased in most of the world due to the increasing importance of non-financial assets. However, it is worth noting that this decline is so low that I tend to disregard it.

Returning to global numbers, the top 10\% richest hold 81\% of the wealth in the world. The distribution of financial and non-financial assets is not uniform. Non-financial assets dominate at the base of the pyramid, gradually being replaced by financial assets as we move up the pyramid. The top 1\% richest hold their wealth primarily in financial assets. In this way, changes in asset prices result in changes in inequality and consequently in average and median wealth.

An impressive fact is that China's growth over 20 years (2000 to 2022) is equivalent to the United States’ growth over 80 years (1925–2005), and the next five years (2022–2027) should be equivalent to an additional 14 years in the U.S., leaving only 8 years of gap. At this pace, China could reach the U.S. by the end of the decade. Another notable point is that during the pandemic, the public sector bailed out the private sector to provide relief by increasing government spending, and in 2022 the altered balance of public versus private wealth began to revert.

